% from main.tex
%%%%%%%%%%  scan idx 527  |  folio 514  %%%%%%%%%%
\[
i_*(b\cdot h^*(a))=i_*(b)\cdot g^*(a),
\]
où $g$ et $h$
sont les morphismes structuraux de $\mathbb P(F)$ et $\mathbb P(E)$. Comme $h=g\circ i$, on a:
\[
i_*(b\cdot h^*(a))=i_*(b\cdot i^*g^*(a))=g^*(a)\cdot i_*(b)
\]
d'après IV 2.12.4.

D'après (3.4.1), on a donc :
\[
K_\bullet(P)\simeq K_\bullet(S)\otimes_{K^\circ(S)}
\bigl(K^\circ(\mathbb P(F))/\operatorname{Im}K^\circ(\mathbb P(E))\bigr).
\]
Soient $\xi_F$ et $\xi_E$ les classes des faisceaux inversibles canoniques sur $\mathbb P(F)$ et $\mathbb P(E)$. Alors, d'après VI 1.1, $K^\circ(\mathbb P(F))$ (resp. $K^\circ(\mathbb P(E))$) admet pour base sur $K^\circ(S)$ les éléments $\xi_F^k$ pour $k=0,\ldots,-n$ (resp. $\xi_E^k$ pour $k=0,\ldots,-n+1$). La formule de projection donne :
\[
i_*(\xi_E^k)=i_*(i^*(\xi_F^k))=\xi_F^k i_*(1).
\]
Si $I$ est l'idéal de $\mathbb P(E)$ dans $\mathbb P(F)$, $I$ est un idéal inversible, et on a $i_*(1)=1-I$. Par ailleurs, $I$ est défini par l'idéal $J$ de $S(F)$ noyau de $S(F)\to S(E)$, qui est engendré par un sous-Module de $F$ isomorphe à $\mathcal O_S$; par suite, $I\simeq \mathcal O_{\mathbb P(F)}(-1)$. L'image de $K^\circ(\mathbb P(E))$ est donc le sous-$K^\circ(S)$-module engendré par les $\xi_F^k-\xi_F^{k-1}$ pour $k=0,\ldots,-n+1$. L'homomorphisme
\[
K^\circ(S)\longrightarrow K^\circ(\mathbb P(F))/\operatorname{Im}(K^\circ(\mathbb P(E)))
\]
est donc un isomorphisme, et par suite, il en est de même pour l'homomorphisme $K_\bullet(S)\to K_\bullet(P)$.

\subsection*{4. Fibrés principaux sous les groupes $\mathrm{GL}(n)_S$}

\textbf{Proposition 4.1.} Soient $S$ un schéma noethérien, $G$ le $S$-groupe $\mathrm{GL}(n)_S$, $P$ un torseur sous $G$, et $E$ le $\mathcal O_S$-Module localement libre de rang $n$ correspondant à $P$. On note $E$ la classe de $E$ dans $K^\circ(S)$, et $J$ le $\lambda$-idéal engendré par $(E-n)$. Si $f$ est le
